% Part: computability
% Chapter: computability-theory
% Section: fixed-point-thm

\documentclass[../../../include/open-logic-section]{subfiles}

\begin{document}

\olfileid[hi]{cmp}{thy}{fix}
\olsection{नियत-बिंदु प्रमेय}

विराम समस्या पर फिर विचार करें। अस्थायी संकेत के रूप में
$\gn{\cfind{x}(y)}$ को $\tuple{x, y}$ के लिए लिखें; इसे मान $\cfind{x}(y)$ के किसी ``नाम''
को निरूपित करने वाला समझें। इस संकेत से विराम समस्या की अनिर्णेयता के
अपने प्रमाणों में से एक को नए शब्दों में कहा जा सकता है।

प्रश्न: क्या निम्न गुण वाला कोई संगणनीय फलन $h$ है? प्रत्येक $x$ और $y$ के लिए,
\[
h(\gn{\cfind{x}(y)}) =
\begin{cases}
1 & \text{यदि $\cfind{x}(y) \fdefined$} \\
0 & \text{अन्यथा।}
\end{cases}
\]

उत्तर: नहीं; अन्यथा आंशिक फलन
\[
g(x) \simeq
\begin{cases}
0 & \text{यदि $h(\gn{\cfind{x}(x)}) = 0$} \\
\text{अपरिभाषित} & \text{अन्यथा}
\end{cases}
\]
संगणनीय होता और इसलिए उसका कोई सूचकांक~$e$ होता। किंतु तब हमें
\[
\cfind{e}(e) \simeq
\begin{cases}
0 & \text{यदि $h(\gn{\cfind{e}(e)}) = 0$} \\
\text{अपरिभाषित} & \text{अन्यथा,}
\end{cases}
\]
मिलता है; तब $\cfind{e}(e)$ तभी और केवल तभी परिभाषित है जब वह परिभाषित
नहीं है—यह विरोधाभास है।

अब $\cfind{e}$ वाले समीकरण को देखें। एक अर्थ में वहाँ स्व-संदर्भ का एक
उदाहरण है: हमने मान $\cfind{e}(e)$ को एक विशेष ढंग से
$\gn{\cfind{e}(e)}$ पर निर्भर बनाया है। नियत-बिंदु प्रमेय कहता है कि
सामान्यतः हम ऐसा {\em कर सकते हैं}---केवल विरोधाभास सिद्ध करने के लिए नहीं।

\olref{lem:fixed-equiv} नियत-बिंदु प्रमेय को कहने के दो समतुल्य ढंग देता
है। तार्किक रूप से कथनों की समतुल्यता उनके दोनों सत्य होने से निकलती है;
पर हमारा वास्तविक आशय है कि प्रत्येक कथन दूसरे से सीधे निष्पन्न होता है,
इसलिए उन्हें उसी प्रमेय के वैकल्पिक कथन माना जा सकता है।

\begin{lem}
\ollabel{lem:fixed-equiv}
निम्नलिखित कथन समतुल्य हैं:
\begin{enumerate}
\item प्रत्येक आंशिक संगणनीय फलन $g(x,y)$ के लिए कोई सूचकांक~$e$ ऐसा है
  कि प्रत्येक~$y$ के लिए
\[
\cfind{e}(y) \simeq g(e,y).
\]
\item प्रत्येक संगणनीय फलन~$f(x)$ के लिए कोई सूचकांक~$e$ ऐसा है कि
  प्रत्येक~$y$ के लिए
\[
\cfind{e}(y) \simeq \cfind{f(e)}(y).
\]
\end{enumerate}
\end{lem}

\begin{proof}
$(1) \Rightarrow (2)$: $f$ मिलने पर $g$ को $g(x,y) \simeq
\fn{Un}(f(x),y)$ से परिभाषित करें। (1) से ऐसा सूचकांक~$e$ पाएँ कि
प्रत्येक~$y$ के लिए
\begin{align*}
\cfind{e}(y) & = \fn{Un}(f(e),y) \\
& = \cfind{f(e)}(y).
\end{align*}

$(2) \Rightarrow (1)$: $g$ मिलने पर $s$-$m$-$n$ प्रमेय से ऐसा $f$ पाएँ
कि प्रत्येक $x$ और~$y$ के लिए $\cfind{f(x)}(y) \simeq g(x,y)$। (2) से
ऐसा सूचकांक~$e$ पाएँ कि
\begin{align*}
\cfind{e}(y) & = \cfind{f(e)}(y) \\
& = g(e,y).
\end{align*}
इससे प्रमाण पूरा होता है।
\end{proof}

\begin{explain}
कथन (1) को सत्य दिखाने (और फलतः (2) को भी) से पहले विचार करें कि वह
कितना विचित्र है। $e$ को कोई कंप्यूटर प्रोग्राम मानें; कथन (1) कहता है
कि कोई भी आंशिक संगणनीय $g(x,y)$ मिलने पर आप ऐसा कंप्यूटर प्रोग्राम $e$
खोज सकते हैं जो $g_e(y) \simeq g(e,y)$ को संगणित करे। दूसरे शब्दों में,
आप ऐसा कंप्यूटर प्रोग्राम खोज सकते हैं जो स्वयं उसी प्रोग्राम का संदर्भ
लेने वाले फलन को संगणित करता है।
\end{explain}

\begin{thm}
\olref{lem:fixed-equiv} के दोनों कथन सत्य हैं। विशेषतः, प्रत्येक आंशिक
संगणनीय फलन $g(x,y)$ के लिए कोई सूचकांक~$e$ ऐसा है कि प्रत्येक~$y$ के लिए
\[
\cfind{e}(y) \simeq g(e,y).
\]
\end{thm}

\begin{proof}
आवश्यक अवयव ऊपर विराम समस्या की चर्चा में पहले से निहित हैं।
$\fn{diag}(x)$ को वह संगणनीय फलन मानें जो प्रत्येक $x$ के लिए फलन
$f_x(y) \simeq \cfind{x}(x,y)$ का कोई सूचकांक लौटाता है, अर्थात्
\[
\cfind{\fn{diag}(x)}(y) \simeq \cfind{x}(x,y).
\]
$\fn{diag}$ को ऐसा फलन मानें जो किसी 2-स्थानी फलन के प्रोग्राम को, मूल
प्रोग्राम को उसके प्रथम तर्क के रूप में नियत करके मिले 1-स्थानी फलन के
प्रोग्राम में बदलता है। फलन $\fn{diag}$ को औपचारिक रूप से इस प्रकार
परिभाषित किया जा सकता है: पहले $s$ को परिभाषित करें:
\[
s(x,y) \simeq \fn{Un}^2(x,x,y),
\]
जहाँ $\fn{Un}^2$ वह 3-स्थानी फलन है जो आंशिक संगणनीय 2-स्थानी फलनों के
लिए सार्वत्रिक है। तब $s$-$m$-$n$ प्रमेय से हम ऐसा आदिम पुनरावर्ती फलन
$\fn{diag}$ खोज सकते हैं जो
\[
\cfind{\fn{diag}(x)}(y) \simeq s(x,y).
\]

अब फलन $l$ को परिभाषित करें:
\[
l(x,y) \simeq g(\fn{diag}(x),y).
\]
और $\gn{l}$ को $l$ का कोई सूचकांक मानें। अंततः
$e = \fn{diag}(\gn{l})$ मानें। तब प्रत्येक $y$ के लिए हमें
\begin{align*}
\cfind{e}(y) & \simeq \cfind{\fn{diag}(\gn{l})}(y) \\
& \simeq \cfind{\gn{l}}(\gn{l}, y) \\
& \simeq l(\gn{l}, y) \\
& \simeq g(\fn{diag}(\gn{l}),y) \\
& \simeq g(e, y),
\end{align*}
मिलता है, जैसा अपेक्षित था।
\end{proof}

\begin{explain}
यहाँ क्या हो रहा है? मान लें कि आपको ऐसा कंप्यूटर प्रोग्राम लिखना है जो
स्वयं को मुद्रित करे। यह भी मान लें कि आप ऐसी प्रोग्रामिंग भाषा में काम
कर रहे हैं जिसमें अक्षर-शृंखला फलनों का समृद्ध और विचित्र संग्रह है।
विशेषतः, मान लें कि आपकी भाषा में कोई फलन $\fn{diag}$ इस प्रकार काम करता
है: आगत अक्षर-शृंखला~$s$ मिलने पर $\fn{diag}$,~$s$ में आए चिह्न `x' के
हर उदाहरण को ढूँढ़ता है और उसे मूल अक्षर-शृंखला के उद्धृत रूप से बदलता
है। उदाहरणतः, अक्षर-शृंखला
\begin{quote}
\begin{verbatim}
hello x world
\end{verbatim}
\end{quote}
को आगत देने पर फलन
\begin{quote}
\begin{verbatim}
hello 'hello x world' world
\end{verbatim}
\end{quote}
निर्गत लौटाता है। उस स्थिति में इच्छित प्रोग्राम लिखना सरल है; आप जाँच
सकते हैं कि
\begin{quote}
\begin{verbatim}
print(diag('print(diag(x))'))
\end{verbatim}
\end{quote}
से कार्य हो जाता है। C++ और Java जैसी अधिक सामान्य प्रोग्रामिंग भाषाओं
में भी यही विचार (अधिक जटिल कार्यान्वयन के साथ) काम करता है।

अब हम नियत-बिंदु प्रमेय के प्रमाण से केवल कुछ चरण दूर हैं। मान लें कि
मुद्रण फलन का कोई रूप $\fn{print}(x,y)$, अक्षर-शृंखला $x$ और दूसरा
संख्यात्मक तर्क $y$ स्वीकार करता है तथा अक्षर-शृंखला $x$ को $y$ बार
मुद्रित करता है। तब ``प्रोग्राम''
\begin{quote}
\begin{verbatim}
getinput(y);
print(diag('getinput(y); print(diag(x), y)'), y)
\end{verbatim}
\end{quote}
आगत $y$ पर स्वयं को $y$ बार मुद्रित करता है।
$\fn{getinput}$---$\fn{print}$---$\fn{diag}$ के ढाँचे को किसी मनमाने
फलन $g(x,y)$ से बदलने पर
\begin{quote}
\begin{verbatim}
g(diag('g(diag(x), y)'), y)
\end{verbatim}
\end{quote}
मिलता है; यह ऐसा प्रोग्राम है जो आगत $y$ पर $g$ को स्वयं उस प्रोग्राम
और $y$ पर चलाता है। ``उद्धृत करने'' को ``के लिए सूचकांक का उपयोग करने'' के रूप
में समझने पर हमें ऊपर का प्रमाण मिलता है।

अभी के लिए यदि आप प्रमाण को औपचारिक बाज़ीगरी अथवा काला जादू मानना चाहें,
तो ठीक है। किंतु आपको ऊपर दिए तर्क का विवरण पुनर्निर्मित कर पाना चाहिए।
जब हम अपूर्णता प्रमेय (और संबंधित ``नियत-बिंदु प्रमेय'') सिद्ध करेंगे,
तब इसके काम करने के कारण को समझने के अन्य ढंगों पर चर्चा करेंगे।
\end{explain}

\begin{tagblock}{lambda}
\begin{digress}
यही विचार ``नियत-बिंदु'' संयोजक प्राप्त करने में प्रयुक्त हो सकता है।
मान लें कि आपके पास लैम्ब्डा पद $g$ है और आप ऐसा दूसरा पद $k$ चाहते हैं
जिसमें $k$, $\beta$-अर्थ में $gk$ के समतुल्य हो। पदों को परिभाषित करें:
\[
\fn{diag}(x) = xx
\]
और
\[
l(x) = g(\fn{diag}(x))
\]
अपने संकेत-संप्रदाय के अनुसार; दूसरे शब्दों में, $l$ पद
$\lambd[x][g(xx)]$ है। $k$ को पद $ll$ मानें। तब हमें
\begin{align*}
k & = (\lambd[x][g(xx)])(\lambd[x][g(xx)]) \\
& \red  g((\lambd[x][g(xx)])(\lambd[x][g(xx)])) \\
& = gk.
\end{align*}
यदि हम
\[
Y = \lambd[g][((\lambd[x][g(xx)])(\lambd[x][g(xx)]))]
\]
लें, तो $Yg$ और $g(Yg)$ एक ही पद में अपचयित होते हैं; अतः
$Yg \equiv_\beta
g(Yg)$। इसे ``करी का संयोजक'' कहते हैं। यदि इसके स्थान पर
\[
Y = (\lambd[xg][g(xxg)])(\lambd[xg][g(xxg)])
\]
लें, तो वास्तव में $Yg$, $g(Yg)$ में अपचयित होता है, जो अधिक प्रबल कथन
है। $Y$ का यह बाद वाला रूप ``ट्यूरिंग का संयोजक'' कहलाता है।
\end{digress}
\end{tagblock}

\end{document}
